% Hướng dẫn cho Tổng biên tập (TBT). Dịch: pdflatex slides.tex (hai lần), sau scripts/test-all.sh
\documentclass[aspectratio=169,12pt]{beamer}
\usepackage{../pi-hd}
\newcommand\vaitro{Hướng dẫn cho Tổng biên tập}
\begin{document}

\bia{Tổng biên tập}{Quyết định về mức và xung đột, duyệt bảng chọn bài, khoá kỳ và đánh số in.}

\chuanbi
\dungthu
\thuatngu

\chu{Việc của Tổng biên tập}{%
\begin{buoc}
\item Đọc mọi bài, phiếu phản biện, thảo luận (kèm email).
\item Quyết định các \textbf{xung đột về mức, tác giả, trùng bài} — chỉ Tổng biên tập đặt „đã giải quyết".
\item Đổi trạng thái khi cần (Không SL, SL-Fail phải ghi lý do); mở lời giải cho phản biện nếu Phụ trách chuyên mục vắng.
\item \textbf{Duyệt} hoặc \textbf{trả lại} bảng chọn bài.
\item \textbf{Khoá kỳ}: đánh số in, gói chế bản cho Ban trình bày.
\end{buoc}}

\manhinh[Đưa về không xoá góp ý, lỗi, người dùng; trước khi đưa về, hệ thống chụp trạng thái hiện tại.]{dung-thu-dua-ve}{Đưa về như trước khi dùng thử}
\manhinh{tbt-xung-dot}{Quyết định xung đột}
\manhinh{tbt-duyet}{Duyệt hoặc trả lại bảng chọn bài}
\manhinh[Khoá kỳ không mở lại được trên trang (đang dùng thử thì Đưa về được): số in ghi vào danh sách đã đăng, bài chuyển trạng thái PL.]{khoa-ky}{Khoá kỳ và đánh số in}

\chu{Khi có sự cố}{%
Thầy cô:
\begin{itemize}\setlength\itemsep{2mm}
\item \textbf{Cần thay một bài sau khi duyệt}: Mở lại để sửa (bấm hai lần để xác nhận) — các bài trở về trạng thái trước khi duyệt.
\item \textbf{Nút Khoá kỳ mờ}: chưa có số in bắt đầu — gõ số, bấm Đánh lại số.
\item \textbf{Bị nhắc „Đánh dấu Tôi đã xem các lưu ý"}: đọc các lưu ý (mục cần kiểm tra, khác mức), đánh dấu rồi khoá.
\item \textbf{Đã khoá nhầm}: không mở lại được trên trang. Đang dùng thử: \emph{Đưa về như trước khi dùng thử}. Ngoài đợt thử:
  báo Quản trị — Quản trị khôi phục từ bản sao lưu đêm trước; việc làm sau bản sao lưu đó phải làm lại.
\end{itemize}}

\gopy

\chu{Tóm tắt}{%
\begin{buoc}
\item Quyết định xung đột ở mục Nguồn và chỉnh sửa của bài.
\item Duyệt bảng (bấm hai lần để xác nhận) — bài thành SL-OK, nhận mức của vị trí; hoặc trả lại có lý do.
\item Khoá kỳ: số in, xem trước tệp .tex, khoá (bấm hai lần để xác nhận).
\end{buoc}}
\end{document}
